% Generated captions for the landscape panels in this SI.
\begin{figure}[p]
  \centering
  \includegraphics[width=0.96\textwidth,height=0.56\textheight,keepaspectratio]{SupplementaryData/figures/landscape_F1.png}
  \caption{F1, Sphere (10D), solver bounds $[-100,100]$, reference optimum $f^*=0$. Panel (a) is the contour view and panel (b) is the 3-D surface of the same slice. The remaining 8 coordinates are fixed at the midpoint of the solver bounds.}
  \label{fig:si_landscape_F1}
\end{figure}
\clearpage

\begin{figure}[p]
  \centering
  \includegraphics[width=0.96\textwidth,height=0.56\textheight,keepaspectratio]{SupplementaryData/figures/landscape_F2.png}
  \caption{F2, Schwefel 2.22 (10D), solver bounds $[-10,10]$, reference optimum $f^*=0$. Panel (a) is the contour view and panel (b) is the 3-D surface of the same slice. The remaining 8 coordinates are fixed at the midpoint of the solver bounds.}
  \label{fig:si_landscape_F2}
\end{figure}
\clearpage

\begin{figure}[p]
  \centering
  \includegraphics[width=0.96\textwidth,height=0.56\textheight,keepaspectratio]{SupplementaryData/figures/landscape_F3.png}
  \caption{F3, Schwefel 1.2 (10D), solver bounds $[-100,100]$, reference optimum $f^*=0$. Panel (a) is the contour view and panel (b) is the 3-D surface of the same slice. The remaining 8 coordinates are fixed at the midpoint of the solver bounds.}
  \label{fig:si_landscape_F3}
\end{figure}
\clearpage

\begin{figure}[p]
  \centering
  \includegraphics[width=0.96\textwidth,height=0.56\textheight,keepaspectratio]{SupplementaryData/figures/landscape_F4.png}
  \caption{F4, Schwefel 2.21 (10D), solver bounds $[-100,100]$, reference optimum $f^*=0$. Panel (a) is the contour view and panel (b) is the 3-D surface of the same slice. The remaining 8 coordinates are fixed at the midpoint of the solver bounds.}
  \label{fig:si_landscape_F4}
\end{figure}
\clearpage

\begin{figure}[p]
  \centering
  \includegraphics[width=0.96\textwidth,height=0.56\textheight,keepaspectratio]{SupplementaryData/figures/landscape_F5.png}
  \caption{F5, Rosenbrock (10D), solver bounds $[-30,30]$, reference optimum $f^*=0$. Panel (a) is the contour view and panel (b) is the 3-D surface of the same slice. The remaining 8 coordinates are fixed at the midpoint of the solver bounds.}
  \label{fig:si_landscape_F5}
\end{figure}
\clearpage

\begin{figure}[p]
  \centering
  \includegraphics[width=0.96\textwidth,height=0.56\textheight,keepaspectratio]{SupplementaryData/figures/landscape_F6.png}
  \caption{F6, Shifted step (10D), solver bounds $[-100,100]$, reference optimum $f^*=0$. Panel (a) is the contour view and panel (b) is the 3-D surface of the same slice. The remaining 8 coordinates are fixed at the midpoint of the solver bounds.}
  \label{fig:si_landscape_F6}
\end{figure}
\clearpage

\begin{figure}[p]
  \centering
  \includegraphics[width=0.96\textwidth,height=0.56\textheight,keepaspectratio]{SupplementaryData/figures/landscape_F7.png}
  \caption{F7, Quartic with noise (10D), solver bounds $[-1.28,1.28]$, reference optimum $f^*=0.5$. Panel (a) is the contour view and panel (b) is the 3-D surface of the same slice. The remaining 8 coordinates are fixed at the midpoint of the solver bounds. The plotted objective is the exact expectation over F7's uniform additive noise.}
  \label{fig:si_landscape_F7}
\end{figure}
\clearpage

\begin{figure}[p]
  \centering
  \includegraphics[width=0.96\textwidth,height=0.56\textheight,keepaspectratio]{SupplementaryData/figures/landscape_F8.png}
  \caption{F8, Schwefel 2.26 (10D), solver bounds $[-500,500]$, reference optimum $f^*=-4189.83$. Panel (a) is the contour view and panel (b) is the 3-D surface of the same slice. The remaining 8 coordinates are fixed at the midpoint of the solver bounds.}
  \label{fig:si_landscape_F8}
\end{figure}
\clearpage

\begin{figure}[p]
  \centering
  \includegraphics[width=0.96\textwidth,height=0.56\textheight,keepaspectratio]{SupplementaryData/figures/landscape_F9.png}
  \caption{F9, Rastrigin (10D), solver bounds $[-5.12,5.12]$, reference optimum $f^*=0$. Panel (a) is the contour view and panel (b) is the 3-D surface of the same slice. The remaining 8 coordinates are fixed at the midpoint of the solver bounds.}
  \label{fig:si_landscape_F9}
\end{figure}
\clearpage

\begin{figure}[p]
  \centering
  \includegraphics[width=0.96\textwidth,height=0.56\textheight,keepaspectratio]{SupplementaryData/figures/landscape_F10.png}
  \caption{F10, Ackley (10D), solver bounds $[-32,32]$, reference optimum $f^*=0$. Panel (a) is the contour view and panel (b) is the 3-D surface of the same slice. The remaining 8 coordinates are fixed at the midpoint of the solver bounds.}
  \label{fig:si_landscape_F10}
\end{figure}
\clearpage

\begin{figure}[p]
  \centering
  \includegraphics[width=0.96\textwidth,height=0.56\textheight,keepaspectratio]{SupplementaryData/figures/landscape_F11.png}
  \caption{F11, Griewank (10D), solver bounds $[-600,600]$, reference optimum $f^*=0$. Panel (a) is the contour view and panel (b) is the 3-D surface of the same slice. The remaining 8 coordinates are fixed at the midpoint of the solver bounds.}
  \label{fig:si_landscape_F11}
\end{figure}
\clearpage

\begin{figure}[p]
  \centering
  \includegraphics[width=0.96\textwidth,height=0.56\textheight,keepaspectratio]{SupplementaryData/figures/landscape_F12.png}
  \caption{F12, Penalized 1 (10D), solver bounds $[-50,50]$, reference optimum $f^*=0$. Panel (a) is the contour view and panel (b) is the 3-D surface of the same slice. The remaining 8 coordinates are fixed at the midpoint of the solver bounds.}
  \label{fig:si_landscape_F12}
\end{figure}
\clearpage

\begin{figure}[p]
  \centering
  \includegraphics[width=0.96\textwidth,height=0.56\textheight,keepaspectratio]{SupplementaryData/figures/landscape_F13.png}
  \caption{F13, Penalized 2 (10D), solver bounds $[-50,50]$, reference optimum $f^*=0$. Panel (a) is the contour view and panel (b) is the 3-D surface of the same slice. The remaining 8 coordinates are fixed at the midpoint of the solver bounds.}
  \label{fig:si_landscape_F13}
\end{figure}
\clearpage

\begin{figure}[p]
  \centering
  \includegraphics[width=0.96\textwidth,height=0.56\textheight,keepaspectratio]{SupplementaryData/figures/landscape_F14.png}
  \caption{F14, Shekel foxholes (2D), solver bounds $[-65.536,65.536]$, reference optimum $f^*=1$. Panel (a) is the contour view and panel (b) is the 3-D surface of the same slice. Both coordinates span the full two-dimensional search domain.}
  \label{fig:si_landscape_F14}
\end{figure}
\clearpage

\begin{figure}[p]
  \centering
  \includegraphics[width=0.96\textwidth,height=0.56\textheight,keepaspectratio]{SupplementaryData/figures/landscape_F15.png}
  \caption{F15, Kowalik (4D), solver bounds $[-5,5]$, reference optimum $f^*=0.0003$. Panel (a) is the contour view and panel (b) is the 3-D surface of the same slice. The remaining 2 coordinates are fixed at the midpoint of the solver bounds.}
  \label{fig:si_landscape_F15}
\end{figure}
\clearpage

\begin{figure}[p]
  \centering
  \includegraphics[width=0.96\textwidth,height=0.56\textheight,keepaspectratio]{SupplementaryData/figures/landscape_F16.png}
  \caption{F16, Six-hump camel (2D), solver bounds $[-5,5]$, reference optimum $f^*=-1.0316$. Panel (a) is the contour view and panel (b) is the 3-D surface of the same slice. Both coordinates span the full two-dimensional search domain.}
  \label{fig:si_landscape_F16}
\end{figure}
\clearpage

\begin{figure}[p]
  \centering
  \includegraphics[width=0.96\textwidth,height=0.56\textheight,keepaspectratio]{SupplementaryData/figures/landscape_F17.png}
  \caption{F17, Branin (2D), solver bounds $[-5,5]$, reference optimum $f^*=0.398$. Panel (a) is the contour view and panel (b) is the 3-D surface of the same slice. Both coordinates span the full two-dimensional search domain.}
  \label{fig:si_landscape_F17}
\end{figure}
\clearpage

\begin{figure}[p]
  \centering
  \includegraphics[width=0.96\textwidth,height=0.56\textheight,keepaspectratio]{SupplementaryData/figures/landscape_F18.png}
  \caption{F18, Goldstein-Price (2D), solver bounds $[-2,2]$, reference optimum $f^*=3$. Panel (a) is the contour view and panel (b) is the 3-D surface of the same slice. Both coordinates span the full two-dimensional search domain.}
  \label{fig:si_landscape_F18}
\end{figure}
\clearpage

\begin{figure}[p]
  \centering
  \includegraphics[width=0.96\textwidth,height=0.56\textheight,keepaspectratio]{SupplementaryData/figures/landscape_F19.png}
  \caption{F19, Hartmann 3-D (3D), solver bounds $[0,1]$, reference optimum $f^*=-3.86$. Panel (a) is the contour view and panel (b) is the 3-D surface of the same slice. The remaining 1 coordinates are fixed at the midpoint of the solver bounds.}
  \label{fig:si_landscape_F19}
\end{figure}
\clearpage

\begin{figure}[p]
  \centering
  \includegraphics[width=0.96\textwidth,height=0.56\textheight,keepaspectratio]{SupplementaryData/figures/landscape_F20.png}
  \caption{F20, Hartmann 6-D (6D), solver bounds $[0,1]$, reference optimum $f^*=-3.32$. Panel (a) is the contour view and panel (b) is the 3-D surface of the same slice. The remaining 4 coordinates are fixed at the midpoint of the solver bounds.}
  \label{fig:si_landscape_F20}
\end{figure}
\clearpage

\begin{figure}[p]
  \centering
  \includegraphics[width=0.96\textwidth,height=0.56\textheight,keepaspectratio]{SupplementaryData/figures/landscape_F21.png}
  \caption{F21, Shekel-5 (4D), solver bounds $[0,10]$, reference optimum $f^*=-10.1532$. Panel (a) is the contour view and panel (b) is the 3-D surface of the same slice. The remaining 2 coordinates are fixed at the midpoint of the solver bounds.}
  \label{fig:si_landscape_F21}
\end{figure}
\clearpage

\begin{figure}[p]
  \centering
  \includegraphics[width=0.96\textwidth,height=0.56\textheight,keepaspectratio]{SupplementaryData/figures/landscape_F22.png}
  \caption{F22, Shekel-7 (4D), solver bounds $[0,10]$, reference optimum $f^*=-10.4029$. Panel (a) is the contour view and panel (b) is the 3-D surface of the same slice. The remaining 2 coordinates are fixed at the midpoint of the solver bounds.}
  \label{fig:si_landscape_F22}
\end{figure}
\clearpage

\begin{figure}[p]
  \centering
  \includegraphics[width=0.96\textwidth,height=0.56\textheight,keepaspectratio]{SupplementaryData/figures/landscape_F23.png}
  \caption{F23, Shekel-10 (4D), solver bounds $[0,10]$, reference optimum $f^*=-10.5364$. Panel (a) is the contour view and panel (b) is the 3-D surface of the same slice. The remaining 2 coordinates are fixed at the midpoint of the solver bounds.}
  \label{fig:si_landscape_F23}
\end{figure}
\clearpage
